% !TEX TS-program = LuaLaTeX

% Author: John H. Lienhard (c) 2026. Reuse under the MIT license: https://ctan.org/license/mit 
% Version 1.06, 2026/07/27

% Documentation: https://ctan.org/pkg/mitthesis

\DocumentMetadata
{
	lang		= en-US,
	pdfstandard = { ua-2, a-4f }, 
	tagging 	= on, 
}

%%%%%%%%%%%%%%%%%%%%%%%%%%%%%%%%%%%%%%%

\documentclass[]{mitthesis}

%% Package for improved typography
\usepackage[nopatch=footnote]{microtype}% typographic fine-tuning (https://ctan.org/pkg/microtype); https://github.com/latex3/tagging-project/issues/67#issue-2163508995


%%%%%%%%%%%  Metadata  %%%%%%%%%%%%%%%%%%%%%%%%%%%%%%%%%%%%%%%%%%%%%%%%%%%%%%%

% Most of the document metadata is created automatically. 
% The following items should be adjusted to match your work. <================= !!!!!!!!!!

\hypersetup{%
	pdfsubject={Template for writing MIT theses with the mitthesis class},
	% Change this to briefly state topic of your thesis 
% 
	pdfkeywords={Massachusetts Institute of Technology, MIT},
	% Add keywords that will help search engines and libraries to find your work.
	% Includes the name[s] of the author[s] 
	% (If you have used \DocumentMetadata, at line 15, you can just put "\CopyrightAuthor," for the names.)
%
	pdfurl={},
	% If you have a url for the thesis, put it here. Otherwise delete this.
	% (MIT Libraries will put your thesis in DSPACE with a persistent url after you submit it.)
%	
	pdfcontactemail={},
	% You can put a [permanent] email address into the metadata, if you like.
	% Otherwise delete this.
%
	pdfauthortitle={},
	% If you have a title, you can include it here.
}

	
%%%%%%%%%%%%%%  End preamble %%%%%%%%%%%%%%%%%%%%%%%%%%%%%%%%%%%%%%%%%%%%%%%%%%%%%%%%%%%%%%%%%%%%%

%%%%%%%%%%%%%%%%%%%%%%%%%%%%%%%%%%%%%%%%%%%%%%%%%%%%%%%%%%%%%%%%%%%%%%%%%%%%%%%%%%%%%%%%%%%%%%%%%%

\begin{document}


%%% edit the following commands to match your thesis %%%%%%%%%%

\title{The Atomic Theory as Applied To Gases, with Some Experiments on the Viscosity of Air}

% \Author{Author full name}{Author department}[Author's first PREVIOUS degree][Author's second PREVIOUS degree][...
% Note that third, fourth, fifth, and sixth arguments are optional [] and may be omitted

% note on names: most of the following names are made up; Silas Holman was a professor at MIT in the 19th century.

\Author{Silas W. Holman}{Department of Nuclear Science and Engineering}%[B.S. Physics, MIT, 1876][M.B.A. Ferengi School of Management, 2022]
%\Author{Luisa Hernández}{Department of Research}%[B.S. Mechanical Engineering, UCLA, 2018][M.S. Stellar Interiors, Vulcan Science Academy, 2020][M.B.A. Ferengi School of Management, 2022]

% For one degree issued by two departments, leave the degree argument blank for the second degree {}.
\Degree{Master of Science in Computational Nuclear Materials}{Department of Nuclear Science and Engineering}
\Degree{}{Center for Computational Science and Engineering}

% If there is more than one supervisor, use the \Supervisor command for each.
\Supervisor{Edward C. Pickering}{Professor of Nuclear Science and Engineering}[Department of Nuclear Science and Engineering]
%\Supervisor{Secunda Castor}{Professor of Research}
%\Supervisor{Quintus Castor}{Professor of Log Dams}

% Professor who formally accepts theses for your department (e.g., the Graduate Officer, Professor Sméagol,...)
% If more than one department, use more than once
% If you need to reduce vertical space, put the acceptor title in the second argument and leave the third blank {}.
\Acceptor{Primus Castor}{Professor and Graduate Officer}{Department of Nuclear Science and Engineering}
\Acceptor{Tertius Castor}{Professor and Graduate Officer}{Center for Computational Science and Engineering}
%\Acceptor{Quarta Castor}{Professor of Lodge Building}{Undergraduate Officer, Department of Mechanical Engineering}

% Usage: \DegreeDate{Month}{year}
% Valid degree months are September, February, May, or June.  
\DegreeDate{June}{1876}

% Date that final thesis is submitted to department
\ThesisDate{May 18, 1876}

% If at least one thesis reader is included, a committee members page will be automatically generated.
% The MIT libraries do not set a format for the committee members page, and if you prefer to use your own format,
% omit ALL readers here and insert your page just before the abstract.
%
\Reader{Marcus Gavius Apicius}{Professor of Cooking Arts}{Department of Food Science}
\Reader{Marie-Antoine Carême}{Professor of Haute Cuisine}{Department of Food Science}
\Reader{Miles Gloriosus}{Professor of Personal Pronouns}{Department of Rhetoric}

%%%%%%  Choose whether to have a CREATIVE COMMONS License  %%%%%%%%%%%%%%%%%%%%%%%%%%%%%%%%%%%%%%
%
% If you are using a cc license, put details of your cc license here. 
% Omit this command if you are not using a cc license.
%
\CClicense{CC BY-NC-ND 4.0}{https://creativecommons.org/licenses/by-nc-nd/4.0/}
%
%%%%%%%%%%%%%%%%%%%%%%%%%%%%%%%%%%%%%%%%%%%%%%%%%%%%%%%%%%%%%%%%%%%%%%%%%%%%%%%%%%%%%%%%%%%%%%%%%

%%% Titlepage
\maketitle

% The abstract environment creates all the required headers and footnote. 
% You only need to add the text of the abstract itself in the file abstract.tex
\begin{abstract}
	\input{abstract.tex}% in this case, use \input rather than \include because you are inside an environment
\end{abstract}

\end{document}	
